\chapter{Iteration 1: Design and Implementation}
\label{ch:iter1}
\section{Iteration Objectives}
Build the application platform and data services that later vision analytics plug into, secure it, and de-risk cross-camera association early with a conservative MVP.

\section{Selected Use Cases / Features}
UC-01, UC-02 to UC-11, UC-12, UC-13 and the operator-facing UC-14 (Chapter~\ref{ch:req}), covering the application/handoff features and the additional FR-21--FR-28 floor-distance prototype.

\section{Design}
The design follows Chapter~\ref{ch:design}: a Next.js frontend with a server-side proxy, a FastAPI service exposing \texttt{/pos/*} and \texttt{/camera/*}, SQLite for POS data, and a file-based handoff output.

\section{Implementation}
The frontend follows the App Router layout, with one folder per page (landing, login, signup, dashboard, customer-info, employee-info, spatial-analytics, sales-trends, reports, alerts, system-status, store-settings, employee-home, my-performance, profile, about). Shared UI primitives (RetroWindow, RetroCard, RetroButton, Navbar, Footer) live in \texttt{src/components/ui}, and helper code such as PDF generation lives in \texttt{src/lib}. The backend lives in \texttt{backend/}.

\subsection{Floor-Distance Command-Line Prototype}
The implementation reuses the existing local trained \path{models/best.pt} weights, whose classes are employee and customer, without retraining or replacement weights. Shared loading/tracking helpers in \path{backend/vision_inference.py} preserve the YOLO/ByteTrack inference route used by the handoff code. The floor CLI selects both role classes; it does not substitute the earlier vertical-position role heuristic.

The entry points are \path{scripts/preview_floor_points.py}, \path{scripts/calibrate_floor.py}, \path{scripts/run_floor_distance.py}, and \path{scripts/evaluate_floor_distance.py}. Geometry, calibration diagnostics, motion monitoring, metric smoothing and evaluation are separate backend modules. The preview runs without calibration; the distance demo exports only supported metre observations and explicit unavailable states. No frontend changes were needed for this prototype.

\section{Completed Features / Module}
\begin{table}[H]
\centering
\caption{Features completed in Iteration 1}
\begin{tabular}{P{0.3\textwidth}P{0.4\textwidth}P{0.17\textwidth}}
\toprule
Feature & Detail & Status\\
\midrule
Manager dashboard pages & Dashboard, Customer Info, Employee Info, Spatial Analytics, Alerts, System Status, Settings & Done (prototype data for vision metrics)\\
Employee workspace & My Workspace and My Performance with rubric and daily log & Done (prototype data)\\
Shared pages & Landing, Login/Signup, About, Profile & Done\\
Reports & Configurable builder, narrative summary, PDF export via jsPDF & Done\\
POS analytics & CSV/XLSX import, product metrics, conversion, forecasts, recommendations, cited assistant & Done (local)\\
Secure access & Signed HttpOnly session cookie and internal secret header & Done (demo credentials)\\
Camera handoff MVP & Offline pipeline, \texttt{/camera/*} jobs, \texttt{/camera-handoff} page & Done (experimental)\\
Floor-distance prototype & Trained roles/track IDs, foot preview, calibration, alignment checks and diagnostic exports & Implemented (local CLI)\\
\bottomrule
\end{tabular}
\end{table}

\section{Technical Implementation Details}
\textbf{Running the system.} \texttt{npm run dev:all} starts Next.js on port 3000 and FastAPI on \texttt{127.0.0.1:8011}. \texttt{npm run pos:api} starts the POS API alone, and \texttt{npm run backend:api} starts the complete backend with \texttt{/pos/*} and \texttt{/camera/*}.

\textbf{Camera handoff.} \path{backend/camera_handoff.py} accepts \texttt{--cam1}, \texttt{--cam2}, and \texttt{--sample-every} arguments and writes \path{backend/output/camera_handoff.json}. The web page \texttt{/camera-handoff} exposes the same batch process through \texttt{/api/camera-handoff}, with configurable normalised boundaries and timing thresholds. Jobs are queued and polled at \texttt{/camera/jobs/\{id\}}.

\textbf{Authentication.} Choosing the Manager role with a valid email and a password of six or more characters creates a server-signed session cookie in the local demo. Manager POS routes require that cookie, so changing \texttt{localStorage} does not unlock them. Secrets are supplied through \texttt{POS\_SESSION\_SECRET} and \texttt{POS\_INTERNAL\_SECRET}.

\subsection{Calibration Configuration and Coordinate Conventions}
Detections use original-frame $[x_1,y_1,x_2,y_2]$ pixels. Ultralytics already restores \texttt{xyxy} boxes after preprocessing, so coordinates are not unletterboxed again. Frame/configuration dimension mismatches, invalid projections and stale homographies are rejected. No lens undistortion is applied; calibration and runtime share the same distorted pixel space. Intrinsic/distortion metadata slots are retained, but correction would require an explicitly consistent update of both paths.

The provisional rectangle follows A--B--D--C: AB = DC = 3.13225 ft and BD = CA = 4.2075 ft, equivalent to width 0.9547098 m and depth 1.282446 m. The configuration records the rectangle geometry and estimated lengths. Dot centres in the 479-by-479 marked screenshot were mapped to the original 1080-by-1080 video using pixel-centre resize coordinates. Frame zero was selected after image-difference comparison and visual inspection; the resized/modified screenshot correspondence remains approximate. Operational coverage is restricted to the marked rectangle.

The rectangle dimensions and original-pixel points are editable in configuration. Regeneration writes the image-to-floor matrix and fitting diagnostics without changing trained weights. Saved measurements are labelled estimated, calibration is provisional and runtime values are approximate. Point-click selection supports numbered correspondences, undo, reset and save; the headless JSON route was exercised.

\subsection{Reproducible Execution}
From the repository root, use the existing virtual environment. The following three command groups reproduce the sampled temporary-video workflow; join wrapped continuation lines when entering each command:
\begin{quote}\small\ttfamily
.venv/Scripts/python.exe scripts/preview\_floor\_points.py\\
\hspace*{1em}--video data/videos/raw/queue\_temp.mp4 --sample-every 30 --headless
\par
.venv/Scripts/python.exe scripts/calibrate\_floor.py\\
\hspace*{1em}--video data/videos/raw/queue\_temp.mp4 --timestamp 0\\
\hspace*{1em}--landmarks data/calibration/queue\_temp.landmarks.json\\
\hspace*{1em}--output data/calibration/queue\_temp.json --headless
\par
.venv/Scripts/python.exe scripts/run\_floor\_distance.py\\
\hspace*{1em}--video data/videos/raw/queue\_temp.mp4\\
\hspace*{1em}--calibration data/calibration/queue\_temp.json\\
\hspace*{1em}--motion-mode monitor --sample-every 30 --headless
\end{quote}
The displayed line breaks are for typesetting; join continuation lines when entering a command. The output directory contains annotated video, top-down video, JSONL observations and summary/coverage records. Actual decoded timestamps are exported. Repository usage instructions are in \path{docs/floor_distance_usage.md}.

\section{Algorithms and Technical Decisions}
\textbf{Handoff algorithm.}
\begin{enumerate}
\item Each camera is processed independently with YOLO detection and local ByteTrack identifiers, sampling every $n$-th frame to reduce CPU cost.
\item Directional boundary events are emitted when a track crosses the configured boundary, which is the right-side exit in cam1 and the left-side entry in cam2.
\item Candidate pairs are filtered by the synchronised time window.
\item Lightweight HSV and shape appearance descriptors are compared for each candidate pair.
\item Each pair is reported as confirmed, uncertain, or unmatched. Confirmed associations receive a stable Master ID for that run.
\end{enumerate}
Direction and appearance are treated as evidence, not proof. Uncertain and unmatched cases are kept in the output on purpose. A verified common cross-camera floor mapping, a third camera, pretrained Re-ID embeddings, timestamp offsets, and real-time guarantees remain future work. The separate single-camera floor homography described below is implemented provisionally.

\textbf{POS analytics.} Metrics are computed deterministically from imported records, so repeated imports give identical results. Assistant answers cite the evidence they use, and forecasts cover 7- and 30-day horizons.

\subsection{Person Floor-Position Approximation}
For a box $B=(x_1,y_1,x_2,y_2)$, the homogeneous floor-contact proxy is
\begin{equation}
\mathbf{p}=\begin{bmatrix}(x_1+x_2)/2 & y_2 & 1\end{bmatrix}^{\mathsf T}.
\label{eq:floor-foot}
\end{equation}
This is an observed proxy rather than proof that feet are visible. Boundary truncation, person overlap and manual exclusion polygons provide provisional rejection signals; seated people and chair occlusion remain limitations. An optional future foot/keypoint locator can replace this step without changing the projection interface.

\subsection{Homography-Based Projection and Distance}
An image-to-floor homography maps the proxy into metre coordinates \cite{hartley2004multiple,opencvHomography}:
\begin{equation}
\mathbf{q}=H_{\mathrm{image}\rightarrow\mathrm{floor}}\mathbf{p},\qquad
X=q_1/q_3,\quad Y=q_2/q_3.
\label{eq:floor-project}
\end{equation}
Non-finite values or a denominator near zero are rejected. For two supported observations from the same camera and timestamp,
\begin{equation}
d_{AB}=\sqrt{(X_A-X_B)^2+(Y_A-Y_B)^2}\quad\text{metres}.
\label{eq:floor-distance}
\end{equation}
Four non-degenerate landmark pairs provide a direct calibration fit. With additional correspondences, the implementation fits floor-to-image using RANSAC with an explicit pixel-space threshold, then inverts the result \cite{fischler1981ransac}. Inliers and fitting residuals are exported as calibration diagnostics.

\subsection{Camera Alignment and Measurement Reliability}
The monitor compares ORB features \cite{rublee2011orb} within the supported floor region, masking detected people and manual exclusions. Matches are checked using RANSAC, inlier count/ratio, spatial spread, residual displacement and supported-region displacement limits. Insufficient evidence produces \texttt{motion\_unverified}; supported alignment shifts produce \texttt{motion\_invalid}. Neither state allows a supported metre output. Modes \texttt{monitor} and \texttt{fixed} both enforce these checks; \texttt{snapshot} supports only the reference frame. Motion compensation is deferred because reliable distributed floor evidence was not established on the temporary clip.

Accepted positions are smoothed in metres using timestamps. Configurable speed/jump gates, short prediction limits, a separate last-reliable-observation time and track expiry preserve quality information. Predictions are excluded from supported pair distances, and alignment failures clear smoothing state. Role labels remain separate observations from the trained detector; their confidence is the joint detection/class score, not two independent probabilities. Selected track pairs or a capped set of all pairs can be exported.

\section{Tools, Frameworks and Technologies}
\begin{table}[H]
\centering
\caption{Technology stack}
\begin{tabular}{P{0.28\textwidth}P{0.6\textwidth}}
\toprule
Layer & Technology\\
\midrule
Frontend & Next.js (App Router), React\\
Styling & Tailwind CSS, custom retro-industrial design system\\
Charts and export & Recharts, jsPDF\\
Computer vision & Trained YOLO role detector (CPU), ByteTrack, OpenCV homography/ORB, NumPy; \texttt{lap} for tracking assignment\\
Agent orchestration & LangGraph\\
POS analytics API & FastAPI, SQLite\\
Fonts & Space Grotesk, Inter, JetBrains Mono\\
Runtime & Node.js 18+, Python 3.12+ (virtual environment)\\
\bottomrule
\end{tabular}
\end{table}

\section{Demonstration Evidence}
\todo{insert screenshots: dashboard, customer info, employee info, reports PDF, sales trends, camera-handoff page, and a sample of camera\_handoff.json}

\subsection{Floor-Distance Demonstration}
Figures~\ref{fig:floor-preview}--\ref{fig:floor-map} show actual local prototype outputs. They demonstrate inference, calibration visualization and export.
\begin{figure}[htbp]
\centering
\includegraphics[width=0.72\textwidth]{ThesisFigs/floor_foot_preview.jpg}
\caption{Trained role detections and track IDs with box-bottom proxies; uncalibrated preview}
\label{fig:floor-preview}
\end{figure}
\begin{figure}[htbp]
\centering
\includegraphics[width=0.72\textwidth]{ThesisFigs/floor_provisional_grid.jpg}
\caption{Provisional A--B--D--C calibration and floor grid based on estimated dimensions}
\label{fig:floor-grid}
\end{figure}
\begin{figure}[htbp]
\centering
\includegraphics[width=0.52\textwidth]{ThesisFigs/floor_top_down.jpg}
\caption{Approximate top-down person-position proxies in the reference snapshot}
\label{fig:floor-map}
\end{figure}

\section{Iteration Outcome}
Iteration 1 delivered a working application platform, a local POS analytics service with deterministic and cited outputs, secure role-aware access, and an experimental camera-handoff pipeline. The vision-derived dashboard metrics still use prototype data, and binding them to real agent output is the first goal of Iteration 2.

The floor-distance CLI prototype additionally demonstrates trained role detection, tracking, provisional calibration and conservative metric export. The quality metadata supports subsequent use of accepted metric observations in downstream interaction analysis.
